% TOP_PROBLEM: {"id": 30006835, "title": "Frey-Mazur Conjecture", "category_id": 1, "category": {"id": 1, "name": "number_theory", "display_name": "Number Theory", "description": "Properties of integers, prime numbers, Diophantine equations.", "slug": "number-theory", "order_index": 1, "created_at": "2026-07-31T15:26:25.670Z"}, "status": "open"}
% FIRST_POSTED: 2026-09-25
% ENTRY_KIND: statement_only
% !TeX program = lualatex
% Definition and sources only; this file is not an LLM solution attempt.
\documentclass[11pt]{article}
\usepackage[margin=25mm]{geometry}
\usepackage{fontspec}
\setmainfont{Latin Modern Roman}
\usepackage{amsmath,amssymb,mathtools}
\usepackage{unicode-math}
\setmathfont{Latin Modern Math}
\usepackage{xurl,hyperref}
\hypersetup{colorlinks=true,urlcolor=blue}
\setcounter{secnumdepth}{0}
\setlength{\parindent}{0pt}
\setlength{\parskip}{5pt}
\setlength{\emergencystretch}{3em}
\begin{document}
\section*{\texorpdfstring{Frey-Mazur Conjecture}{Frey-Mazur Conjecture}}\label{30006835}
\subsection{Definitions and mathematical statement}
For an elliptic curve over ${\mathbb{Q}}$, $E[p]$ denotes the two-dimensional ${\mathbb{F}}_p$-space of geometric $p$-torsion, with its action of $G_{{\mathbb{Q}}}=\operatorname{Gal}(\overline{{\mathbb{Q}}}/{\mathbb{Q}})$. A ${\mathbb{Q}}$-isogeny is a nonzero morphism of elliptic curves over ${\mathbb{Q}}$ preserving the origin, necessarily with finite kernel. The selected Frey--Mazur assertion is
\[
 p>17,\qquad E[p]\cong E'[p]\text{ as }G_{{\mathbb{Q}}}\text{-modules}
 \quad\Longrightarrow\quad E\sim_{{\mathbb{Q}}}E'.
\]
The isomorphism is not additionally required to preserve the Weil pairing. The explicit cutoff $17$ is part of the target; existence of some unspecified larger universal cutoff is weaker.
\par\smallskip\textit{Formulation and concept references:} \href{https://annals.math.princeton.edu/2016/184-3/p02}{[S1]}.

\subsection{Short English statement}
For every prime greater than seventeen, does identical Galois-module structure on two rational elliptic curves' p-torsion force the curves to be isogenous over the rationals?
\subsection{Sources}
\begin{itemize}
\item[S1]Benjamin Bakker and Jacob Tsimerman, "p-torsion monodromy representations of elliptic curves over geometric function fields," Annals of Mathematics 184 (2016), 709-744.
\url{https://annals.math.princeton.edu/2016/184-3/p02}.
\textit{Formulation source.}
\end{itemize}
\end{document}
